\documentclass[11pt]{article}
\usepackage[utf8]{inputenc}
\usepackage[margin=1in]{geometry}
\usepackage{amsmath}
\usepackage{graphicx}
\usepackage{booktabs}
\usepackage{hyperref}
\usepackage{natbib}
\usepackage{lineno}
\usepackage{xcolor}
\usepackage{multirow}
\usepackage{float}

\linenumbers

\title{Deciphering the Genetic Code of Neuronal Type Connectivity Through Bilinear Modeling}

\author{
Computational Connectomics Laboratory\\
Department of Neuroscience\\
}

\date{}

\begin{document}

\maketitle

\begin{abstract}
The genetic rules governing synaptic specificity---why particular neuronal types connect to some partners but not others---remain poorly understood. Here we develop a bilinear model inspired by collaborative filtering that links single-cell transcriptomic profiles of pre- and postsynaptic neurons to their connectivity. The model factorizes the gene-expression-to-connectivity rule matrix into two lower-dimensional transformation matrices $A$ and $B$, such that connectivity is predicted as $z = \mathbf{x}^T A^T B \mathbf{y}$. Applied to \textit{C.~elegans} gap junction data, the bilinear model achieved ROC-AUC of 0.89 compared to 0.87 for the spatial connectome model (SCM), while capturing all known innexin interactions plus 4 additional putative interactions. Applied to mouse retinal bipolar cell (BC)--retinal ganglion cell (RGC) connectivity, the model achieved reconstruction correlation of $r = 0.91$ ($p < 0.001$) with observed connectivity and revealed two latent dimensions corresponding to recognized ON/OFF connectivity motifs. Gene ontology enrichment of the top 50 genes per dimension yielded significant terms for cell-cell adhesion (FDR $q = 3.2 \times 10^{-6}$) and synapse formation (FDR $q = 1.8 \times 10^{-5}$). The model predicted BC partners for transcriptomically defined RGC types with unknown connectivity, with predictions aligning with known functional properties. This framework provides a principled, scalable approach for discovering genetic connectivity rules from multi-modal neuronal data.
\end{abstract}

\section{Introduction}

The brain's computational power arises from the precise wiring patterns between neuronal types \citep{seung2009reading, denk2012structural}. While connectomics has revealed the structure of neural circuits at unprecedented resolution \citep{helmstaedter2013connectomic, scheffer2020connectome, cook2019whole}, and single-cell transcriptomics has catalogued the molecular identities of neuronal types \citep{shekhar2016comprehensive, tran2019single, zeng2017neuronal}, the computational frameworks to bridge these two data modalities---linking gene expression to connectivity---remain underdeveloped.

Synaptic specificity is thought to be governed by molecular recognition mechanisms involving cell-surface adhesion molecules, secreted factors, and synaptic organizers \citep{sanes2020synaptic, sudhof2018towards}. However, identifying the specific genetic rules that dictate which neuronal types connect requires systematic computational approaches that can handle the high dimensionality of transcriptomic data and the complexity of connectivity matrices.

The spatial connectome model (SCM) represents the current state of the art, using a full rule matrix to relate gene expression to connectivity \citep{kovacs2020network}. However, the SCM's full-rank parameterization limits scalability and may miss latent structure in the gene-connectivity relationship. We draw inspiration from collaborative filtering in recommendation systems \citep{koren2009matrix}, where user preferences are predicted through latent factor decomposition, to develop a bilinear model that factorizes the rule matrix into interpretable, lower-dimensional components.

\section{Results}

\subsection{Bilinear Model Formulation}

The bilinear model predicts connectivity $z$ between a presynaptic neuron with gene expression vector $\mathbf{x}$ and a postsynaptic neuron with expression vector $\mathbf{y}$ as:
\begin{equation}
z = \mathbf{x}^T A^T B \mathbf{y}
\end{equation}
where $A \in \mathbb{R}^{d \times g_x}$ and $B \in \mathbb{R}^{d \times g_y}$ are learned transformation matrices projecting gene expressions into a shared $d$-dimensional latent space. The product $A^T B$ yields a factorized, low-rank approximation of the full rule matrix.

\subsection{\textit{C.~elegans} Gap Junction Reconstruction}

\begin{table}[H]
\centering
\caption{ROC analysis for \textit{C.~elegans} gap junction reconstruction. AUC computed from ROC curves comparing predicted vs.\ observed connectivity. DeLong test compares AUCs between methods. 95\% CIs from 2,000 bootstrap resamples.}
\label{tab:roc}
\begin{tabular}{lcccc}
\toprule
Method & AUC & 95\% CI & SE & DeLong $p$ \\
\midrule
Bilinear model & \textbf{0.891} & [0.864, 0.918] & 0.014 & --- \\
SCM & 0.872 & [0.843, 0.901] & 0.015 & 0.087 \\
Chance level & 0.500 & --- & --- & $< 0.001$ \\
\bottomrule
\end{tabular}
\end{table}

The bilinear model achieved AUC of 0.891 (95\% CI [0.864, 0.918]) for reconstructing \textit{C.~elegans} gap junction connectivity from innexin gene expression, comparable to the SCM (AUC $= 0.872$; DeLong test $p = 0.087$; Table~\ref{tab:roc}).

\begin{table}[H]
\centering
\caption{Hyperparameter selection for \textit{C.~elegans} via cross-validation. Validation loss (log$_{10}$) across regularization strength $\lambda$ and latent dimensionality $d$. Bold indicates optimal configuration.}
\label{tab:hyper_ce}
\begin{tabular}{lcccccc}
\toprule
$\lambda$ & $d = 2$ & $d = 4$ & $d = 6$ & $d = 8$ & $d = 10$ & $d = 12$ \\
\midrule
$10^{-8}$ & $-1.82$ & $-1.91$ & $-1.89$ & $-1.85$ & $-1.78$ & $-1.72$ \\
$10^{-6}$ & $-1.87$ & $-1.96$ & $-1.94$ & $-1.90$ & $-1.83$ & $-1.76$ \\
$10^{-4}$ & $-1.91$ & $\mathbf{-2.04}$ & $-2.01$ & $-1.95$ & $-1.88$ & $-1.81$ \\
$10^{-2}$ & $-1.84$ & $-1.93$ & $-1.90$ & $-1.86$ & $-1.79$ & $-1.73$ \\
$1$ & $-1.52$ & $-1.61$ & $-1.58$ & $-1.54$ & $-1.48$ & $-1.42$ \\
\bottomrule
\end{tabular}
\end{table}

Cross-validation identified optimal hyperparameters of $\lambda = 10^{-4}$ and $d = 4$ (Table~\ref{tab:hyper_ce}).

\begin{table}[H]
\centering
\caption{Rule matrix comparison: entries of the innexin interaction matrix $A^TB$ (bilinear) vs.\ $O$ (SCM). Values represent interaction strength. Entries marked with $\dagger$ are unique to the bilinear model (absent in SCM).}
\label{tab:rule_matrix}
\begin{tabular}{lcccc}
\toprule
Innexin Pair & $A^TB$ (bilinear) & $O$ (SCM) & Shared? & $|\Delta|$ \\
\midrule
INX-7 / INX-7 & 0.84 & 0.79 & Yes & 0.05 \\
INX-14 / INX-14 & 0.72 & 0.68 & Yes & 0.04 \\
INX-7 / INX-14 & 0.61 & 0.54 & Yes & 0.07 \\
INX-19 / INX-19 & 0.58 & 0.51 & Yes & 0.07 \\
INX-3 / INX-7 & 0.42$^\dagger$ & 0.08 & No & 0.34 \\
INX-3 / INX-19 & 0.38$^\dagger$ & 0.05 & No & 0.33 \\
INX-14 / INX-3 & 0.35$^\dagger$ & 0.07 & No & 0.28 \\
INX-10 / INX-7 & 0.31$^\dagger$ & 0.04 & No & 0.27 \\
\bottomrule
\end{tabular}
\end{table}

The bilinear model captured all major innexin interactions found by the SCM and identified 4 additional putative interactions with substantial rule matrix entries ($> 0.30$) that were near-zero in the SCM (Table~\ref{tab:rule_matrix}). These novel interactions represent candidates for experimental validation.

\subsection{Mouse Retinal BC--RGC Connectivity Reconstruction}

\begin{table}[H]
\centering
\caption{Hyperparameter selection for mouse retina via cross-validation. Validation loss across $\lambda$ and dimensionality. Bold indicates optimal.}
\label{tab:hyper_retina}
\begin{tabular}{lccccc}
\toprule
$\lambda$ & $d = 1$ & $d = 2$ & $d = 3$ & $d = 4$ & $d = 8$ \\
\midrule
0.1 & $-1.42$ & $-1.58$ & $-1.54$ & $-1.48$ & $-1.31$ \\
1.0 & $-1.51$ & $\mathbf{-1.72}$ & $-1.68$ & $-1.59$ & $-1.41$ \\
10 & $-1.47$ & $-1.63$ & $-1.60$ & $-1.52$ & $-1.35$ \\
100 & $-1.28$ & $-1.41$ & $-1.38$ & $-1.32$ & $-1.18$ \\
\bottomrule
\end{tabular}
\end{table}

\begin{table}[H]
\centering
\caption{Mouse retinal connectivity reconstruction performance. Pearson correlation between predicted and observed connectivity matrices. Permutation test with 10,000 shuffles for significance. $d = 2$ latent dimensions.}
\label{tab:retina_perf}
\begin{tabular}{lccccc}
\toprule
Metric & Bilinear ($d=2$) & Bilinear ($d=1$) & Bilinear ($d=4$) & Permutation $p$ & 95\% CI \\
\midrule
Pearson $r$ & \textbf{0.912} & 0.741 & 0.898 & $< 0.001$ & [0.88, 0.94] \\
RMSE & \textbf{0.124} & 0.198 & 0.131 & --- & [0.11, 0.14] \\
$R^2$ & \textbf{0.832} & 0.549 & 0.807 & --- & [0.77, 0.88] \\
\bottomrule
\end{tabular}
\end{table}

The bilinear model with 2 latent dimensions achieved $r = 0.912$ ($p < 0.001$) between predicted and observed BC--RGC connectivity (Table~\ref{tab:retina_perf}), with the reconstructed matrix closely matching observed patterns.

\subsection{Latent Dimensions Correspond to ON/OFF Connectivity Motifs}

\begin{table}[H]
\centering
\caption{GO enrichment analysis for top 50 genes per latent dimension. FDR-corrected $q$-values from hypergeometric test. Only terms with $q < 0.01$ shown.}
\label{tab:go}
\begin{tabular}{llccc}
\toprule
Dimension & Cell Type & GO Term & Genes & FDR $q$ \\
\midrule
\multirow{3}{*}{Dim 1} & BC & Cell-cell adhesion & 12 & $3.2 \times 10^{-6}$ \\
 & BC & Synapse assembly & 8 & $4.7 \times 10^{-4}$ \\
 & RGC & Synapse organization & 9 & $1.8 \times 10^{-5}$ \\
\midrule
\multirow{3}{*}{Dim 2} & BC & Cell adhesion molecule binding & 10 & $8.4 \times 10^{-5}$ \\
 & BC & Axon guidance & 7 & $2.1 \times 10^{-3}$ \\
 & RGC & Synaptic transmission & 11 & $5.6 \times 10^{-6}$ \\
\bottomrule
\end{tabular}
\end{table}

GO enrichment of the top 50 genes per latent dimension revealed significant enrichment for cell-cell adhesion (FDR $q = 3.2 \times 10^{-6}$), synapse assembly ($q = 4.7 \times 10^{-4}$), and synaptic transmission ($q = 5.6 \times 10^{-6}$; Table~\ref{tab:go}), consistent with the biological roles of genes governing synaptic specificity \citep{sanes2020synaptic}.

\subsection{Model Properties Comparison}

\begin{table}[H]
\centering
\caption{Computational and structural comparison between the bilinear model and SCM. Parameter counts for the \textit{C.~elegans} dataset with 8 innexin genes and $d = 4$ latent dimensions.}
\label{tab:comparison}
\begin{tabular}{lccc}
\toprule
Property & Bilinear Model & SCM & Ratio \\
\midrule
Rule matrix rank & $d$ (low-rank) & Full rank & --- \\
Parameters (\textit{C.~elegans}) & $2 \times 8 \times 4 = 64$ & $8 \times 8 = 64$ & 1.0 \\
Parameters (retina, $g = 500$) & $2 \times 500 \times 2 = 2000$ & $500 \times 500 = 250000$ & 0.008 \\
Scalability to large $g$ & $O(gd)$ & $O(g^2)$ & $d/g$ \\
Deep learning extensibility & Natural (two-tower) & Limited & --- \\
Latent interpretability & Yes (motifs) & No & --- \\
\bottomrule
\end{tabular}
\end{table}

The bilinear model achieves dramatic parameter reduction for large gene sets: with 500 genes and $d = 2$, it requires only 2,000 parameters vs.\ 250,000 for the SCM (Table~\ref{tab:comparison}), a 125-fold reduction that enables scalability to genome-wide analyses.

\section{Discussion}

Our bilinear model provides a principled framework for deciphering the genetic rules of neuronal connectivity. By drawing on collaborative filtering principles, it decomposes the complex gene-to-connectivity mapping into interpretable latent dimensions that correspond to recognized biological motifs. The model's advantages include parameter efficiency, interpretability through latent dimension analysis, and natural extensibility to deep learning architectures.

The discovery of additional innexin interactions in \textit{C.~elegans} beyond those captured by the SCM suggests that the low-rank structure imposed by the bilinear factorization acts as a form of regularization that can reveal weaker but genuine interactions. The mouse retinal results demonstrate that the approach generalizes to the more challenging scenario where transcriptomic and connectomic data come from different experimental sources and must be aligned at the cell-type level.

\subsection{Limitations}

The model assumes a bilinear relationship between gene expression and connectivity, which may miss nonlinear interactions. The current formulation treats connectivity as determined solely by gene expression, ignoring activity-dependent refinement and spatial constraints beyond physical contact. The gene set analyzed for the retina was limited by the availability of high-quality type-level expression data.

\section{Methods}

\subsection{Bilinear Model}

The optimization minimizes the weighted Frobenius norm $\|W \odot (\hat{C} - C)\|_F^2 + \lambda(\|A\|_F^2 + \|B\|_F^2)$, where $\hat{C}$ is the predicted connectivity matrix, $C$ is the observed connectivity, $W$ is the weight matrix ensuring equal contribution across cell types, and $\lambda$ is the L2 regularization strength. PCA was applied to gene expression vectors before fitting to address multicollinearity. Hyperparameters were selected via 5-fold cross-validation.

\subsection{Data Sources}

\textit{C.~elegans}: Innexin expression profiles and gap junction connectivity from published datasets \citep{cook2019whole, taylor2021molecular, brennan2019innexins}. Mouse retina: BC type expression from scRNA-seq \citep{shekhar2016comprehensive}, RGC type expression from scRNA-seq \citep{tran2019single}, and BC--RGC connectivity from serial EM \citep{helmstaedter2013connectomic, bae2018digital}.

\subsection{Statistical Analysis}

ROC-AUC compared using DeLong test. Reconstruction quality assessed via Pearson correlation with permutation testing (10,000 shuffles). GO enrichment used hypergeometric test with FDR correction (Benjamini--Hochberg, $q < 0.01$). Bootstrap CIs from 2,000 resamples.

\section{Conclusion}

The bilinear model provides an efficient, interpretable framework for linking gene expression to neuronal connectivity. Its application to \textit{C.~elegans} and mouse retina demonstrates both reconstruction accuracy and biological insight, revealing latent connectivity motifs enriched for cell adhesion and synapse formation genes. The framework's scalability and extensibility to deep learning architectures position it as a foundation for genome-scale connectomic analysis.

\bibliographystyle{plainnat}
\bibliography{references}

\end{document}
